%%%%%%%%%%%%%%%%%%%%%%%%%%%%%%%%%%%%%%%%%%%%%%%%%%%%%%%%%%%%%%%%%%%%%%%%%%%%%%%%%%%
% Plantilla para os traballos TFG/TFM da EPS de Enxeñaría
% LaTeX Template
% Version 1.1 (12/6/18)
%
% Autor: Xosé R. Fdez-Vidal
% Subdirección de Calidade da EPS de Enxeñaría
%
% Esta plantilla está baseada no template de:
% Steve Gunn (http://users.ecs.soton.ac.uk/srg/softwaretools/document/templates/)
%
% Template license:
% CC BY-NC-SA 3.0 (http://creativecommons.org/licenses/by-nc-sa/3.0/)
%%%%%%%%%%%%%%%%%%%%%%%%%%%%%%%%%%%%%%%%%%%%%%%%%%%%%%%%%%%%%%%%%%%%%%%%%%%%%%%%%%%%%%%

%----------------------------------------------------------------------------------------
%	PAQUETES E OUTRAS CONFIGURACIÓNS DO DOCUMENTO
%----------------------------------------------------------------------------------------

\documentclass[
11pt, % Tamaño da fonte do documento, opcións: 10pt, 11pt, 12pt
%oneside, % A dobre cara (alternando os marxes), descomenta para cambiar a un documento por unha cara soamente
spanish, % Idioma. Opcións: galician, spanish, english, ..
singlespacing, % Interliñado simple, alternativas: onehalfspacing ou doublespacing
%draft, % Descomentar para facilitar o modo draft (sen figuras, sen links, ...)
%nolistspacing, % Seo documento é onehalfspacing ou doublespacing, descomenta esta opción para asignar o espaciado nas listas a simple
%liststotoc, % Descomenta para engadir a lista de figuras/tablas/etc á táboa de contidos
%toctotoc, % Descomenta para engadir a táboa principal de contidos á táboa de contidos
parskip, % Descomenta para engadir un espazo maior entre parágrafos
%nohyperref, % Descomenta para non cargar o paquete hyperref (ligazóns no documento)
headsepline, % Descomenta para engadir unha liña debaixo da cabeceira
chapterinoneline, % Descomenta para poñer o título do capítulo a continuación do número da páxina
%consistentlayout, % Descomenta para cambiar o layout da declaración, abstract e agradecementos e que tome os valores da clase por defecto
]{./cls/TFG_TFM_EPS} % O ficheiro da clase que especifica o estilo do documento


%----------------------------------------------------------------------------------------
%	ESTABLECEMENTO DO XESTOR DE BIBLIOGRAFIA, O ESTILO E O FICHEIRO QUE CONTEN AS REFERENCIAS
%----------------------------------------------------------------------------------------

%\usepackage[backend=bibtex,style=authoryear,natbib=true]{biblatex} % Empregamos bibtex para o manexo da bibliografía con estilo de referencia autor-ano (paracido a norma APA)
\usepackage[backend=bibtex,style=ieee,natbib=true]{biblatex} %Estilo IEEE
%tes mais estilos, só tes que buscar as opcións existentes no paquete biblatex na opción style

\addbibresource{bibliografia/exemplo.bib} % Nome do ficheiro bib onde temos a bibliografía
\usepackage[autostyle=true]{csquotes} % Requirido para xerar as comillas na bibliografía en función do idioma


%Ollo:
\graphicspath{{Figuras/}{Figuras/cap1/}{Figuras/cap2/}{Figuras/cap3/}{./}} % ESPECIFICA OS PATHS ONDE ESTAN AS FIGURAS. Aumenta os paht tanto como capitulos teñas ou cartafoles crees no directorio de bibliografia


%----------------------------------------------------------------------------------------
%	ASIGNACIÓN DE MARXES (esto podes cambialo a teu gusto)
%----------------------------------------------------------------------------------------

\geometry{
	paper=a4paper, % Tamaño do folio A4
	inner=2.5cm, % Marxe interior
	outer=3.8cm, % Marxe exterior
	bindingoffset=.5cm, % Binding offset
	top=1.5cm, % Marxe superior
	bottom=1.5cm, % Marxe inferior
	%showframe, % Descomenta para visualizar os cadros que representan o diseño feito da páxina
    }


%----------------------------------------------------------------------------------------

% Definimos algúns comandos para mater o formatodo independente do contido
\newcommand{\keyword}[1]{\textbf{#1}}
\newcommand{\tabhead}[1]{\textbf{#1}}
\newcommand{\code}[1]{\texttt{#1}}
\newcommand{\file}[1]{\texttt{\bfseries#1}}
\newcommand{\option}[1]{\texttt{\itshape#1}}

%----------------------------------------------------------------------------------------


%----------------------------------------------------------------------------------------
%	INFORMACIÓN DO TFG/TFM (DÉBEO COMPLETAR CADA ALUMNO ASPIRANTE)
%----------------------------------------------------------------------------------------

%Datos referentes aos titores do traballo (comenta os que non precises)
\supervisor{\href{}{Dr. \textsc{Xosé Ramón Fernández Vidal}}} % Nome do titor 1, usado no paxina do título, amosase co comando \supname
\department{\href{}{Dpto. de Física Aplicada}} % Nome do departamento ao que pertence o titor académico e amósase co comando \deptname

\supervisorseg{Dr. \textsc{Marcos Fernández González}} % Nome do titor 2, usado no paxina do título, amosase co comando \supnameseg
\departmentseg{\href{}{Dpto. de Electrónica y Computación}} % Nome do departamento ao que pertence o titor académico e amósase co comando \deptnameseg

% \supervisorter{\textsc{Juan Motel}} % Nome do titor 2, usado no paxina do título, amosase co comando \supnameter
% \departmentter{\href{https://gominolassl.com}{Gominolas SL}} % Nome do departamento ao que pertence o titor académico e amósase co comando \deptnameter



%Datos referentes ao alumno e ao seu traballo
\thesistitle{\textsc{Análisis comparativa de los métodos de segmentación 2D y 3D aplicados a imágenes médicas}} % Título do traballo, móstrase co comando \ttitle
\author{\textsc{Alfredo Vázquez Rodríguez}} % O teu nome, imprimese na páxina do título e na do resume,  amosase co comando \authorname
\def \dniautor {58.00.24.76 H} %dni do autor
\keywords{\itshape{Electrónica, Aprendizaxe máquina, (mínimo 5 palabras!!)}} % Keywords do teu traballo, amosanse na páxina do resumo do traballo, aparecen onde poñas o comando \keywordnames
\def \dataConvocatoria {Julio 2026} %Data da convocatoria. Requirida por normativa!


%----------------------------------------------------------------------------------------
%
%----------------------------------------------------------------------------------------

%Metadatos no pdf xerado
\AtBeginDocument{
\hypersetup{pdftitle=\ttitle} % Mete no pdf o título do teu traballo
\hypersetup{pdfauthor=\authorname} % Asigna no pdf o teu nome como autor do traballo
\hypersetup{pdfkeywords=\keywordnames} % Mete no pdf as túas keywords
}

%Empezamos o documento!!!
\begin{document}

\nocite{*}
\frontmatter % Empregamos números romanos para numerar as páxinas pre-contidos no indice (i, ii, iii, iv...)

\pagestyle{plain} % Establecemos o estilo da cabeceira ao defecto ata que no corpo do documento se asigne o estilo definido no preambulo.

%----------------------------------------------------------------------------------------
%	PORTADA (TAPA)
%----------------------------------------------------------------------------------------

\begin{titlepage}
	\newpage
	\hfill
	\begin{textblock*}{\paperwidth}(6mm,260mm)
		\includegraphics[width=12cm]{LogoEPS}
	\end{textblock*}

	\begin{textblock*}{\paperwidth}(130mm,0mm)
		\colorbox{azulUSC}{\makebox[0.38\paperwidth][r]{\shortstack[r]{\vspace{77.5mm}}}}
	\end{textblock*}


	\begin{textblock*}{\paperwidth}(0mm,80mm)
		\colorbox{verdeCTERRA}{\makebox[\paperwidth][c]{\begin{minipage}[c][4cm]{0.9\paperwidth}{\begin{center}
			\huge\rmfamily\color{white}{\textbf{\ttitle}}
		\end{center}}\end{minipage} }}
	\end{textblock*}

	\begin{textblock*}{5cm}(130mm,122.6mm)
		\colorbox{azulUSC}{\hspace{5mm} \parbox{0.35\paperwidth}{
			\hypersetup{hidelinks}
			\vspace{4cm}
			\color{white}
			\textbf{\textsc{\tipotraballo}}
			\\\textbf{\textsc{ Para la titulación de}}
			\\\textbf{\textsc{\degreename}}
			\vskip 1cm
			\textbf{\textsc{Autor}:
			\\ \authorname
			\vskip 1cm
			\textsc{Tutores}:
			\\\supname % Nome do Titor 1
			\\{\small \deptname}
			\vskip 0.2cm
			\supnameseg   % Nome do Titor 2
			\\{\small \deptnameseg}
			\vskip 0.2cm
			\textsc{Convocatoria}:
			\\ \dataConvocatoria
			\vfill
			\vskip 10cm
		}}}
	\end{textblock*}
	\begin{textblock*}{\paperwidth}(6mm,8mm)
		\includegraphics[width=5.5cm]{LogoTerra}
	\end{textblock*}
\end{titlepage}



%----------------------------------------------------------------------------------------
%	PÁXINA DE TÍTLO
%----------------------------------------------------------------------------------------

\begin{titlepage}
\begin{center}

%\vspace*{.01\textheight}
{\scshape\LARGE \univname\par}\vspace{1.5cm} % NOME da USC
\textsc{\Large \tipotraballo}\\[0.5cm] % Tipo de Traballo (TFG/TFM)

\HRule \\[0.4cm] % liña horizontal
{\huge \bfseries \ttitle\par}\vspace{0.4cm} % Titulo do traballo
\HRule \\[1.5cm] % liña horizontal

\begin{minipage}[t]{0.4\textwidth}
\begin{flushleft} \large
\textsc{Autor:}\\
\href{}{\authorname} %Nome do autor - elimina o que hai dentro das chaves do comando \href para eliminar a ligazón
\end{flushleft}
\end{minipage}
\begin{minipage}[t]{0.4\textwidth}
\begin{flushright} \large
	\textsc{Tutores:} \\
	\supname % Nome do Titor 1
	\\{\small \deptname}
	\vskip 0.2cm
	\supnameseg   % Nome do Titor 2
	\\{\small \deptnameseg}

\end{flushright}
\end{minipage}\\[2cm]

\vfill

\large \textit{Este trabajo se presenta para cumplir con los requirimentos normativos\\ esigidos para lograr el \degreename}\\[0.3cm] % Texto requirido
\textit{en la}\\[0.3cm]

\large \textit{\facname.}\\[0.4cm]

%\groupname\\\deptname\\[2cm] % Grupo de investigación e departamento

\vfill

{\large Fecha de la Convocatoria: \dataConvocatoria}\\[4cm] % Data

\begin{textblock*}{\paperwidth}(8mm,260mm)
	\includegraphics{LogoEPS} % Logo do EPS
\end{textblock*}
\vfill
\end{center}
\end{titlepage}

%----------------------------------------------------------------------------------------
%	AUTORIZACIÓN BiBLIOTECA
%----------------------------------------------------------------------------------------

\begin{declaration}
\addchaptertocentry{Autorización} % Engade a declaracion á táboa de contidos
\noindent
Yo, \textsc{\authorname}, alumno de la titulación {\degreename} de la \textit{\facname} de la \textit{\univname}, con DNI {\dniautor}, autorizo la ubicación de la seguinte copia de mí \textbf{Trabajo Fin de Grado} en la biblioteca de Intercentros para que pueda ser consultada por las personas que lo deseen.


\vspace{5em}

\noindent Firmado:\\
\rule[0.5em]{25em}{0.5pt} % This prints a line for the signature
\vspace{2em}

\noindent Fecha:\\
\nopagebreak
\parbox[b]{25em}{ % Creamos una caja del mismo ancho que la línea
    \centering \today \\ % Centramos la fecha
    \rule[0.5em]{25em}{0.5pt} % Dibujamos la línea justo debajo
}
\vfill
\end{declaration}

\cleardoublepage


%----------------------------------------------------------------------------------------
%	PÁXINA PARA UNHA CITA QUE CHE GUSTE
%----------------------------------------------------------------------------------------

\vspace*{0.2\textheight}
\noindent{\itshape Tenta non converterte nunha persoa de éxito, senón nunha persoa de valor.}\bigbreak
\hfill Albert Einstein





%----------------------------------------------------------------------------------------
%	DESEÑO DA PÁXINA QUE CONTÉN O RESUMO DO TRABALLO
%----------------------------------------------------------------------------------------
%%%%%%%%%%%%%%%%%%%%%%%%%%%%%%%%%%%%%%
%Resumo en galego
%%%%%%%%%%%%%%%%%%%%%%%%%%%%%%%%%%%%%%
\renewcommand{\abstractname}{Resumo}
\renewcommand{\byname}{por}
\begin{abstract}[\textsc{Análise comparativa dos métodos de segmentación 2D e 3D aplicados a imaxes médicas}]
	\addchaptertocentry{\abstractname} % Engade o resumo a táboa de contidos
	\lipsum[1]

\end{abstract}



%%%%%%%%%%%%%%%%%%%%%%%%%%%%%%%%%%%%%%
%Resumo en castelan
%%%%%%%%%%%%%%%%%%%%%%%%%%%%%%%%%%%%%%
\renewcommand{\abstractname}{Resumen}
\renewcommand{\byname}{por}
\renewcommand{\palabrasChave}{Palabras Clave: }
\begin{abstract}
	\addchaptertocentry{\abstractname} % Engade o resumo a táboa de contidos
	\lipsum[1]

\end{abstract}



%%%%%%%%%%%%%%%%%%%%%%%%%%%%%%%%%%%%%%
%Resumo en inglés
%%%%%%%%%%%%%%%%%%%%%%%%%%%%%%%%%%%%%%
\renewcommand{\abstractname}{Abstract}
\renewcommand{\byname}{by}
\renewcommand{\palabrasChave}{Keywords: }
\begin{abstract}[\textsc{Comparative Analysis of 2D and 3D Segmentation Methods in Medical Imaging}]
	\addchaptertocentry{\abstractname} % Engade o resumo a táboa de contidos
	\lipsum[1]

\end{abstract}



%----------------------------------------------------------------------------------------
%	AGRADECEMENTOS
%----------------------------------------------------------------------------------------

\begin{acknowledgements}
\addchaptertocentry{\acknowledgementname} % Metemos agradecementos no índice
\lipsum[1]
\end{acknowledgements}

%----------------------------------------------------------------------------------------
%	LISTA DE CONTIDOS/FIGURAS/TÁBOAS
%----------------------------------------------------------------------------------------
\renewcommand{\contentsname}{Índice General} %para renomear en galego o indice
\tableofcontents % inserta a táboa principal de contidos (índice xeral)

%\renewcommand{\listfigurename}{Índice de Figuras} %para renomear en galego o indice
\listoffigures % inserta o índice de figuras


%Cambiamos os nomes ao galego

\renewcommand{\listtablename}{Índice de Tablas} %para renomear en galego o indice
\listoftables % inserta o índice de táboas

%----------------------------------------------------------------------------------------
%	ABREVIACIÓNS
%----------------------------------------------------------------------------------------

\begin{abbreviations}{ll} % Inclúe unha lista de abreviacion que empregues no traballo (unha táboa de dúas columnas)


\end{abbreviations}

%----------------------------------------------------------------------------------------
%	CONSTANTES/OUTRAS DEFINICIÓNS
%----------------------------------------------------------------------------------------

\begin{constants}{lr@{${}={}$}l} % Lista de constantes nunha táboa de 3 columnas

% O comando \SI{}{} está no paquete siunitx do sistema intenacioanl de unidades, Mira a documentación para empregalo correctamente http://osl.ugr.es/CTAN/macros/latex/contrib/siunitx/siunitx.pdf
%Este paquete da a ortotipografía matematica para os símbolos matemáticos e as unidades correspondetes do SI. Empregao!!


\end{constants}

%----------------------------------------------------------------------------------------
%	SIMBOLOS
%----------------------------------------------------------------------------------------

\begin{symbols}{lll} % Lista de símbolos nunha táboa de 3 columnas



\end{symbols}

%----------------------------------------------------------------------------------------
%	DEDICATORIA
%----------------------------------------------------------------------------------------

\dedicatory{Adicado a miña familia \ldots}

%----------------------------------------------------------------------------------------
%	CONTIDO DO TRABALLO - CAPÍTULOS
%----------------------------------------------------------------------------------------

\mainmatter % Empezamos a numeración das páxinas con números normais (1,2,3...)


\pagestyle{thesis} % Retornamos ao formato de tese

% Incluímos os capítulos do traballo como ficheiros separados e situados no cartafol Capitulos
% Descomenta as liñas a midida que avances na redación dos capítulos
% Capítulo 1

\chapter{Estrutura do traballo}

\label{Chapter1} % para referenciar este capítulo en calquera parte deste documento, emprega \ref{Chapter1} 



\section{Preliminares}
Recoméndase consultar a norma UNE 50136:1997: documentación, presentación de tese e documentos similares. Tamén existen recursos en liña que axudan a organizar todo o proceso de elaboración do Traballo. 

Ten en conta, que a morfoloxía do documento non está condicionada pola plantilla e podes estruturalo acorde ao tipo de documento que vaias a realizar (proxecto construtivo, traballo técnico ou de investigación).

\subsection{Cuberta}
Pode utilizarse como plantilla a cuberta que se inclúe neste asegurándose de que se inclúen os datos precisos. Debe conter todos os datos que se especifican: nome completo do autor, título do traballo, nome completo do titor, data, título ao que se opta, etc.
\subsection{Portada}
A cuberta é a tapa do documento e é a primeira folla que aparece ao abrir a tapa do documento. Igualmente pode utilizarse como modelo a que se inclúe aquí e debe conter todos os datos que figuran nesta: o nome da Escola, a titulación e, no seu caso, a intensificación, o título do TFG/TFM, os nomes do autor, do titor e, no seu caso, do relator, a área de coñecemento, o departamento e o ano de execución do traballo.

\subsection{Resumo}
Texto que mostra de forma abreviada e precisa o contido do traballo. Deberase incorporar unha versión en galego, castelán e inglés que se colocarán en orde consecutivo. 
\subsection{Prefacio ou Introdución}
Debe incluír unha breve explicación das razóns que levaron á realización do Traballo, o propósito e os obxectivos que se pretenden, o ámbito, alcance e límites da investigación, así como a metodoloxía empregada e, se se considera oportuno, un avance das conclusións alcanzadas. 
\subsection{Índice xeral}
Táboa de contidos onde se reflicten todas as partes do Traballo e os seus anexos, se os houbese. Deben aparecer os títulos, na súa orde e con indicación da páxina na que se poden atopar.
Se o Traballo consta de varios volumes, cada un deberá levar o seu propio índice, pero débese incluír tamén un índice xeral.
\subsection{Lista de ilustracións e táboas}
Se o Traballo inclúe ilustracións e táboas pódese engadir unha listaxe que inclúa o número identificativo que figura dentro do texto, a lenda e o número da páxina na que se atopa.
Tamén é conveniente mencionar os datos sobre as fontes de onde se obtiveron ditas ilustracións se non se incluíron no propio texto da memoria.
\subsection{Lista de abreviaturas e símbolos}
O Traballo debe conter as abreviaturas e símbolos internacionalmente recoñecidos. Se se incorporan unidades, abreviaturas ou acrónimos que poidan ser pouco coñecidos deberanse explicar brevemente nestas listas.
\subsection{Glosario}
Os termos que requiran definición ou explicación deberanse incorporar nun glosario.

\section{Texto principal}
Para a súa elaboración é importante ter en conta a súa distribución en capítulos e seccións numeradas. 
O texto debe comezar cunha introdución que mostre as investigacións previas existentes sobre o tema e destacar os obxectivos e métodos seguidos para levar a cabo a investigación ou análise do tema tratado.
Para finalizar débense escribir as conclusións que deben estar en relación cos obxectivos marcados previamente.

\section{Bibliografía}
Debe conter as referencias bibliográficas dos documentos consultados para demostrar as bases do traballo realizado e avalar os datos incorporados e citados no texto.
Elaborarase de forma normalizada, para o que se aconsella utilizar a norma UNE vixente (actualmente a UNE ISO 690:2013. Información e documentación. Directrices para a redacción de referencias bibliográficas e de citas de recursos de información).
Para a elaboración desta parte do Traballo recoméndase consultar a Web, guías e axudas para a elaboración das referencias bibliográficas.

\section{Anexos}
Pódese incluír desta forma material extenso utilizado no traballo, importante para xustificar os resultados e as conclusións obtidas, pero que non é esencial para a comprensión do texto principal. Poden ser datos estatísticos, lexislación, etc.
A paxinación debe ser correlativa e continuar a do texto principal. Cada un dos anexos debe identificarse cunha letra maiúscula do alfabeto, comezando pola letra A, precedida da palabra Anexo.

% Capítulo 2

\chapter{Breve guía de uso deste patrón \LaTeX{}}

\label{Chapter2} % para referenciar este capítulo en calquera parte deste documento, emprega \ref{Chapter1} 


\section{Benvidos e grazas}
Benvidos a este patrón \LaTeX{} para traballos fin de grao/máster, que pretende facilitar o inicio dos alumnos para escribir o seu traballo en \LaTeX{}.

Se estas escribindo o teu traballo (ou o farás no futuro próximo) que é técnico e ten ecuacións matemáticas, entón facelo en \LaTeX{} é altamente recomendable e ademais, a forma de facelo máis segura e sen preocuparte polo formatado e pola ortotipografía matemática correcta. Estes aspectos, son os que levan moito tempo en procesadores de texto convencionais e, case sempre, non cumpren coas regras ISO establecidas. 

En \LaTeX{} é moi fácil escribir documentos profesionais e moi longos en calquera ordenador sen especiais requirimentos de memoria RAM. Funciona como un linguaxe de marcas, mediante comandos simples, que automaticamente formatan os índices de contidos, marxes, cabeceiras e pés de páxina de xeito consistente e sen esforzo. Unha das principais fortalezas é a gran calidade na ortotipografía das ecuacións matemáticas e a súa rápida escritura en texto.

%----------------------------------------------------------------------------------------

\section{Aprendendo \LaTeX{}}

\LaTeX{} non é un procesador de textos \textsc{wysiwyg} (What You See is What You Get) como pode ser Microsoft Word ou Apple's Pages. Polo contrario, un documento escrito para \LaTeX{} é un ficheiro de texto puro que non contén ningún tipo \emph{formato}. Un dálle a \LaTeX{} as instrucións de como debe formatar o documento final escribindo os comandos dentro do documento, por exemplo, se queres empregar letra itálica  como \emph{texto en itálica}, escribes o comando \verb|\emph{text}| e pos o texto que desexes entre as chaves. Entón \LaTeX{}, cando atope este comando, poñerá ese texto en cursiva. Neste aspecto, \LaTeX{} é moi semellante ao linguaxe HTML.

\subsection{Unha (non tan corta) Introdución a \LaTeX{}}

Se eres un neófito en \LaTeX{}, hai un bó libro -- libre on-line PDF -- chamado, \enquote{La introducción no-tan-corta a \LaTeX{}}. Podes baixar a última versión dende a URL: \url{http://mirrors.ctan.org/info/lshort/spanish/lshort-a4.pdf}


É recomendable adicar un pouco de tempo a aprender o básico de \LaTeX{} e crear varios documentos básicos para ter un mínimo adestramento. Podes empregar patróns de distintos tipos de documentos na ena URL: \url{http://www.LaTeXTemplates.com}\\ 


\subsection{Símbolos matemáticos comúns para \LaTeX{}}
Hai multitude de símbolos matemáticos dispoñibles para \LaTeX{} e non é posible lembrarse de todos os comandos para cada un deles. A maioría dos que ti podes empregar, podes consultar os seus comandos nesta URL: \url{http://www.sunilpatel.co.uk/latex-type/latex-math-symbols/}. Nesta páxina hai un listado co símbolo matemático en grande e debaixo del o comando \LaTeX{} correspondente.

%----------------------------------------------------------------------------------------

\section{Empezando a escribir o teu traballo a partir deste patrón}

Se xa estas familiarizado con \LaTeX{}, entón deberías explorar a estrutura de cartafoles e ficheiros deste patrón e proceder a colocar a información do teu \emph{TRABALLO TFG/TFM} nos bloques do ficheiro \file{main.tex}. A sección \ref{FillingFile} na páxina \pageref{FillingFile} axudaranche neste proceso. Asegúrate de ler a sección \ref{ThesisConventions} acerca das convencións empregades neste patrón (para traballos de investigación porque para proxectos construtivos ou de outro tipo as conveccións varían e debes consultalas co teu titor).


Antes de empezar a utilizar este patrón deberías asegurarte que eres capaz de compilar este documento. Se queres modificar algún aspecto das páxinas iniciais, debes retocar o ficheiro de estilo \file{TFG\_TFM\_EPS.cls} que está dentro do cartafol chamado \file{cls}.


%----------------------------------------------------------------------------------------

\section{Que inclúe este patrón?}

\subsection{Cartafoles}
Este patrón está empaquetado nun ficheiro zip que ao descomprimilo expande unha estrutura de cartafoles e ficheiros. Os nomes dos cartafoles son autoexplicativos: 

\keyword{Apendices} -- Neste cartafol colocarás os apéndices. Cada apéndice debería ter o seu propio ficheiro \file{.tex} separado. Un exemplo e unha patrón base de apéndice van incluídos neste cartafol.

\keyword{Capitulos} -- Neste cartafol é onde debes colocar os ficheiros do teu TFG/TFM. Un traballo pode ter un número variable de capítulos dependendo da natureza do TFG/TFG (investigación, proxecto, etc.). Cada capítulo debería ter o seu propio ficheiro \file{.tex} e, como aproximación para un traballo de investigación, poden ser divididos en:
\begin{itemize}
\item Capítulo 1: Introdución ao tema do TFG/TFM 
\item Capítulo 2: Revisión bibliográfica
\item Capítulo 3: Modelo teórico/Metodoloxía experimental (laboratorio, etc.)
\item Capítulo 4: Detalles do experimento 1
\item Capítulo 5: Detalles do experimento 2
\item Capítulo 6: Discusión dos resultados experimentais
\item Capítulo 7: Conclusións e dirección futuras
\end{itemize}
Esta estrutura de capítulos é a máis habitual en ciencias experimentais, pero podes adaptala a natureza do teu TFG/TFM (Apéndices, Anexos, etc.).

\keyword{Figuras} -- Este cartafol contén todas as figuras que empregues na memoria e están estruturadas en subcartafoles para cada capítulo ou apéndice para unha mellor organización. A medida que engadas novos subcartafoles inclúe a ruta na instrución \keyword{graphicspath} do \file{main.tex}.
Os ficheiros das figuras poden ser calquera formato: png, tif, jpg, pdf, etc. pero recomendamos empregar formatos vectoriais para unha mellor calidade, como son o pdf ou eps.


\subsection{Ficheiros}

Tamén se inclúen varios ficheiros, a maioría deles son de texto puro, e podes ver o seu contido con calquera editor de texto. Despois da compilación inicial, veras que se crean varios ficheiros auxiliares polo \LaTeX{} ou BibTeX dos cales non precisas preocuparte e podes eliminalos sen efectos secundarios:

\keyword{exemplo.bib} -- Este é un ficheiro importante que contén toda a información bibliográfica e referencias que ti citarás no teu documento TFG/TFM empregando BibTeX. Podes escribilo manualmente os seus items pero hai xestores de citas bibliográficas que che permiten exportar as referencias ao formato BibTeX o cal que simplifica a cantidade de traballo que tes que facer.

\keyword{TFG\_TFM\_EPS.cls} -- Este é un ficheiro importante. Este é o ficheiro de clase que lle di a \LaTeX{} como é formato o documento. 

\keyword{main.pdf} -- Este é o documento final do teu traballo TFG/TFM (en formato PDF) xerado por \LaTeX{} e listo para imprimir.

\keyword{main.tex} -- Este é outro ficheiro importante. Nel lle dicimos a \LaTeX{} que ficheiros se deben incluír na compilación. Contén a estrutura do documento e está comentado para facilitar o entendemento do que está a facer instrución.

Ficheiros que \emph{non} están incluídos, pero que son xerados por \LaTeX{} como ficheiros auxiliares:

\keyword{main.aux} -- ficheiro auxiliar xerado por \LaTeX{}, se se elimina, o \LaTeX{} o rexenera cando volvas a recompilar.

\keyword{main.bbl} -- ficheiro auxiliar xerado por BibTeX. No ficheiro \file{.bib} contén todas as referencias bibliográficas, o ficheiro \file{.bbl} contén as referencias que citas no teu traballo e empregase para construír a sección de bibliografía no TFG/TFM.

O resto de ficheiros auxiliares xerados, pero que non explicaremos (porque carecen de importancia para un usuario normal), son: \keyword{main.blg}, \keyword{main.lof}, \keyword{main.log}, \keyword{main.lot} e \keyword{main.out} 

Desta longa lista, só os ficheiros con extensión \file{.bib}, \file{.cls} and \file{.tex} son os máis importantes. Os ficheiros auxiliares poden ser ignorados ou eliminados posto que \LaTeX{} e BibTeX os rexenerarán cada vez que recompiles o documento principal.

%----------------------------------------------------------------------------------------

\section{Introducindo a información no \file{main.tex}}\label{FillingFile}

Precísase que personalices a información referente ao teu traballo, título, titor, tribunal, etc. Edita co TexStudio (ou calquera outro editor de \LaTeX{}) o \file{main.tex} no bloque 1 etiquetado como \emph{INFORMACIÓN DO TFG/TFM}.

Cando completes isto, salva o ficheiro e recompila \code{main.tex}. Toda esta información aportada se colocará automaticamente nas páxinas da portada, agradecementos, etc.


%----------------------------------------------------------------------------------------

\section{Explicación do ficheiro \code{main.tex}}\label{ThesisConventions}

O ficheiro \file{main.tex} contén a estrutura da tese e está documentado mediante comentarios escritos que explican que páxinas, seccións e formato esta creando o código \LaTeX{}. Cada elemento deste documento está dividido en bloques comentados con títulos en letras maiúsculas que indican o contido deste bloques.
	
As páxinas co índice de contidos, lista de figuras e táboas están pensadas para que se creen automaticamente. Non obstante, para certos documentos técnicos precisase definir unha lista de abreviacións empregadas e/ou constantes e témolo que facer a man. Non obstante, se non precisas isto no teu traballo, podes comentalas para que non aparezan no teu documento final. 
	
A lista de símbolos é dividida segundo os alfabetos romano e grego, mentres que os símbolos das abreviacións está pensado que sexa un listado alfabético (obviamente, isto tes que facelo ti manualmente porque LaTeX{} non o pode facer por ti).
	

Finalmente, hai o bloque onde os capítulos son incluídos. Descomenta as liñas (eliminar o \code{\%} do código) cando escribes os capítulos. Cada capítulo debería ser escrito no seu ficheiro propio e posto no cartafol \emph{Capítulos}, nomeados como \file{Capítulo1}, \file{Capítulo2},\ldots O mesmo procedemento serve para os apéndices, cada un tería que ir no seu ficheiro propio e colocado no \emph{cartafol de Apéndices}.
		
Respecto a bibliografía, o estilo de empregado (opción \option{authoryear} pero podes mudalo) inclúe os autores e o ano entre parénteses. Adicionalmente, inserta unha ligazón á cita o cal facilita moito ao lector atopar a referencia bibliográfica de interese. 

%----------------------------------------------------------------------------------------

\subsection{Formato impreso}

Esta patrón está deseñado para obter un documento para ser impreso a dobre cara como corresponde a maioría deste tipo de documentos. Podemos cambiar isto e imprimir o documento por unha soa cara simplemente descomentando a opción \option{oneside} no comando \code{documentclass} situado ao inicio do ficheiro \file{main.tex}. 

A cabeceira das páxinas pares conteñen o número de páxina na parte esquerda e o título da capítulo na parte dereita. Nas páxinas impares, teñen na parte interior o título da sección e na exterior a numeración das páxinas. 
 
As fontes son asignadas a 11 puntos por defecto e o interliñado do texto é simple, non obstante, podes cambiar estas opcións ao inicio do ficheiro \file{main.tex} cambiando os valores de \option{singlespacing} por \option{onehalfspacing} ou \option{doublespacing}.



\subsection{Referencias}

O paquete \code{biblatex} empregase para formatar a bibliografía e as referencias insírense no texto coia instrución co comando \parencite{Reference1}. As opcións empregadas no ficheiro \file{main.tex} referéncianse cos nomes de autores e data. Múltiples referencias sepáranse con puntos e comas (e.g. \parencite{Reference2, Reference1}) e as referencias con máis de tres autores só se amosa o primeiro seguido \emph{et al.} indicando que hai máis autores (e.g. \parencite{Reference3}). Esto faise automaticamente. Para ver como usar as referencias, bótalle un ollo ao código fonte do ficheiro \file{Capitulo1.tex} e visita esta URL: \url{https://www.latex-tutorial.com/tutorials/bibtex/}.

A bibliografía organizase alfabeticamente polo primeiro apelido do primeiro autor se elixes o estilo de referencias APA, pero se elixes o estilo IEEE, ordenaranse polo orden no que foron citados no texto. Agora vou citar varios tipos de documentos (artigos, conferencias, webs, libros, etc) para que vexas o resultado na sección de bibliografía (estilo IEEE) \parencite{Reference2, conference, inbook, incollection, book, manual, mastersthesis, misc, phdthesis, proceedings, techreport, unpublished, web}.


\subsection{Cadros}

As cadros ou táboas son unha forma importante de visualizar os resultados, abaixo, é un exemplo de táboa que é xerada por este código:

{\small
\begin{verbatim}
\begin{table}
\caption{Os efectos do tratamento X e Y sobre catro grupos de estudo.}
\label{tab:treatments}
\centering
\begin{tabular}{l l l}
\toprule
\tabhead{Grupos} & \tabhead{Tratamento X} & \tabhead{Tratamento Y} \\
\midrule
1 & 0.2 & 0.8\\
2 & 0.17 & 0.7\\
3 & 0.24 & 0.75\\
4 & 0.68 & 0.3\\
\bottomrule\\
\end{tabular}
\end{table}
\end{verbatim}
}

\begin{table}
\caption{Os efectos do tratamento X e Y sobre catro grupos de estudo.}
\label{tab:treatments}
\centering
\begin{tabular}{l l l}
\toprule
\tabhead{Grupos} & \tabhead{Tratamento X} & \tabhead{Tratamento Y} \\
\midrule
1 & 0.2 & 0.8\\
2 & 0.17 & 0.7\\
3 & 0.24 & 0.75\\
4 & 0.68 & 0.3\\
\bottomrule\\
\end{tabular}
\end{table}

Podes referencias as táboas co comando \verb|\ref{<label>}| onde a etiqueta (label) defínese dentro do entorno da táboa. Mira \file{Capitulo1.tex} para un exemplo de etiqueta e citación (e.g. Táboa~\ref{tab:treatments}).

Unha ferramenta on-line e interactiva que xera código \LaTeX de táboas podes atopala na URL \url{http://www.tablesgenerator.com/latex_tables}.

\subsection{Figuras}

A maneira de inserir figuras na memoria é empregar o código como o seguinte:
\begin{verbatim}
\begin{figure}
\centering
\includegraphics[width=5cm]{Figures/Electron}
\decoRule
\caption[Un Electrón]{Un electrón (ilustración artistica).}
\label{fig:Electron}
\end{figure}
\end{verbatim}
Poñendo este código produce a figura do que electróns que estas vendo debaixo. 

\begin{figure}[th]
\centering
\includegraphics[width=4cm]{Electron}
%\decoRule
\caption[Un Electrón]{Un electrón (ilustración artística).}
\label{fig:Electron}
\end{figure}

Algunhas veces as figuras non aparecen situadas no texto onde ti colocas as instrucións anteriores. A situación depende do espazo que queda na páxina para colocar a figura. Se non hai suficiente sitio, \LaTeX{} colocará a ilustración ao principio da seguinte páxina. Non obstante, sempre toma as mellores decisións e ti non tes porque preocuparte por iso (tamén pode forzar a \LaTeX a facer o que ti desexas!).


As figuras deben ter pés de figuras (captions) e se precisas referirte a el (como na  Figure~\ref{fig:Electron}). O comando \verb|\caption| contén dúas partes, a primeira, dentro dos corchetes está o título que aparecerá na \emph{Lista de Figuras}, e debe ser breve. Na segunda parte,dentro das chaves deberá conter o texto máis longo e explicativo que aparecerá no pé da figura.

O comando \verb|\decoRule| é opcional e simplemente pon unha barra horizontal para mellorar a estética. Se fas isto para unha imaxe, faino en todo o texto para manter a homoxeneidade.

\LaTeX{} é capaz de ler imaxes en formatos pdf, jpg e png.

\subsection{Escribindo fórmulas}

Se o teu traballo ten contido matemático amplo, ten por seguro que \LaTeX{} é a mellor elección para este tipo de contidos. Non obstante, mira a referencia dada para aprender o mínimo requirido de \LaTeX en \enquote{Unha Introdución no-tan-breve a \LaTeX} ou esta guía matemática \enquote{A Short Math Guide to \LaTeX} que podes descargar dende: \url{ftp://ftp.ams.org/pub/tex/doc/amsmath/short-math-guide.pdf}


Hai diferentes símbolos en \LaTeX{} que hai que lembrar, afortunadamente podes atopar os mais comúns en \href{http://ctan.org/pkg/comprehensive}{The Comprehensive \LaTeX~Symbol List}.

Podes escribir unha ecuación, que é automaticamente numerada por \LaTeX{}, con este código:
\begin{verbatim}
\begin{equation}
E = mc^{2},
\label{eqn:Einstein}
\end{equation}
\end{verbatim}

que producirá a ecuación de Einstein máis famosa:
\begin{equation}
E = mc^{2},
\label{eqn:Einstein}
\end{equation}

Todas as ecuacións que escribes (que non están medio de dous parágrafos) son automaticamente numeradas por \LaTeX{}. Se non queres que se numeren, emprega a forma:
\begin{verbatim}
\[ a^{2}=4. \]
\end{verbatim}

Para as ecuacións ou símbolos matemáticas que van na liñas de texto empregase $a^{2}=4$ os símbolos de dólar: \begin{verbatim} $ a^{2}=4 $ \end{verbatim}



Nos primeiros intentos de escribir fórmulas en \LaTeX, podes empregar a ferramenta codecogs (\url{http://www.codecogs.com/latex/eqneditor.php}), similar ao editor de ecuacións do Word, que  xera os comandos \LaTeX da fórmula seleccionada.

%----------------------------------------------------------------------------------------

\section{Sección e subseccións}

Debes estruturar o TFG en capítulos, seccións e subseccións. \LaTeX{} constrúe o índice de contidos automaticamente mirando os títulos dentro das chaves dos comandos \verb|\chapter{}|, \verb|\section{}|  e \verb|\subsection{}| que escribes no teu documento.

O índice de contidos só listará tres niveis. Un \verb|chapter{}| é o nivel (0). Un \verb|\section{}| é o nivel  (1) e \verb|\subsection{}| é o nivel (2). No teu documento tamén poderás empregar o comando \verb|subsubsection{}|, o cal será o nivel (3). A profundidade do índice de contidos podes configuralo no ficheiro \file{TFG\_TFM\_EPS.cls}.
%----------------------------------------------------------------------------------------
\vspace{5em}

Boa sorte co teu TFG/TFM en \LaTeX!!
\begin{flushright}
patróns para distintos tipos de documentos \LaTeX{} libres---\\
\href{http://www.LaTeXTemplates.com}{LaTeXTemplates.com}
\end{flushright}

%\include{Capitulos/Capitulo3}
%\include{Capitulos/Capitulo4}
%\include{Capitulos/Capitulo5}


%----------------------------------------------------------------------------------------
%	CONTIDO DO TRABALLO - APÉNDICES
%----------------------------------------------------------------------------------------

\appendix % Comando para decirlle a LaTeX que os seguintes capítulos son Apéndices

% Incluímos os apéndices do traballo como ficheiros separados e situados no cartafol Apendices
% Descomenta as liñas a midida que avances na redación dos Apendices

% Apendice A

\chapter{FAQs} % Titulo

\label{ApendiceA} % Para referenciar este apéndice en calquera parte do documento, emprega \ref{ApendiceA}

\section{Como cambiar as cores das ligazóns?}

As cores das ligazóns no documento poden ser cambiadas como se indica de seguido:

{\small\verb!\hypersetup{urlcolor=red}!}, ou

{\small\verb!\hypersetup{citecolor=green}!}, ou

{\small\verb!\hypersetup{allcolor=blue}!}.

\noindent Se queres ocultar as ligazóns no documento, emprega:

{\small\verb!\hypersetup{allcolors=.}!}, ou mellor: 

{\small\verb!\hypersetup{hidelinks}!}.

\noindent Se queres ter as ligazóns no pdf pero non no texto impreso, usa:

{\small\verb!\hypersetup{colorlinks=false}!}.

%% Apendice A

\chapter{FAQs} % Titulo

\label{ApendiceA} % Para referenciar este apéndice en calquera parte do documento, emprega \ref{ApendiceA}

\section{Como cambiar as cores das ligazóns?}

As cores das ligazóns no documento poden ser cambiadas como se indica de seguido:

{\small\verb!\hypersetup{urlcolor=red}!}, ou

{\small\verb!\hypersetup{citecolor=green}!}, ou

{\small\verb!\hypersetup{allcolor=blue}!}.

\noindent Se queres ocultar as ligazóns no documento, emprega:

{\small\verb!\hypersetup{allcolors=.}!}, ou mellor: 

{\small\verb!\hypersetup{hidelinks}!}.

\noindent Se queres ter as ligazóns no pdf pero non no texto impreso, usa:

{\small\verb!\hypersetup{colorlinks=false}!}.

%% Apendice A

\chapter{FAQs} % Titulo

\label{ApendiceA} % Para referenciar este apéndice en calquera parte do documento, emprega \ref{ApendiceA}

\section{Como cambiar as cores das ligazóns?}

As cores das ligazóns no documento poden ser cambiadas como se indica de seguido:

{\small\verb!\hypersetup{urlcolor=red}!}, ou

{\small\verb!\hypersetup{citecolor=green}!}, ou

{\small\verb!\hypersetup{allcolor=blue}!}.

\noindent Se queres ocultar as ligazóns no documento, emprega:

{\small\verb!\hypersetup{allcolors=.}!}, ou mellor: 

{\small\verb!\hypersetup{hidelinks}!}.

\noindent Se queres ter as ligazóns no pdf pero non no texto impreso, usa:

{\small\verb!\hypersetup{colorlinks=false}!}.


%Para incluir Anejos en formato pdf (saída dun programa, etc) ao teu documento os podes metelo dentro do ficheiro dun apéndice empregando o seguinte comando (dentro do ficheiro .tex do apéndice)
%Nota: para incluir documento pdf emprega este comando:
%\includepdf[pages=inicial-final]{nome_do_documento}




%----------------------------------------------------------------------------------------
%	BIBLIOGRAFÍA
%----------------------------------------------------------------------------------------

\printbibliography[heading=bibintoc]

%----------------------------------------------------------------------------------------

\end{document}
